\documentclass[../main.tex]{subfiles}
\begin{document}
本章主要参考\cite{mcs}。

% Reference: https://doi.org/10.1587/transfun.2022EAL2095
本章汇集离散数学中的逻辑、证明方法、关系运算、格与布尔代数。集合、映射、等价关系与偏序的基础定义见第\ref{chap:set-theory}章；计数、递推、生成函数及图论见第\ref{chap:combinatorics-and-graph-theory}章；整除、同余及初等数论公式见第\ref{chap:elementary-number-theory}章。

\section{命题逻辑}
\subsection{联结词与真值表}\index{联结词与真值表|hyperpage}
命题具有确定的真值，以 $1$ 表示真、$0$ 表示假。否定、合取、析取、蕴含与等价分别记为 $\neg,\land,\lor,\to,\leftrightarrow$；异或记为 $\oplus$。析取允许两者同时为真。
\[
\begin{array}{cc|c|ccccc}
p&q&\neg p&p\land q&p\lor q&p\to q&p\leftrightarrow q&p\oplus q\\\hline
0&0&1&0&0&1&1&0\\
0&1&1&0&1&1&0&1\\
1&0&0&0&1&0&0&1\\
1&1&0&1&1&1&1&0
\end{array}
\]
蕴含仅在前件真、后件假时为假；$p\to q$ 不表示因果关系。$p$ 是 $q$ 的充分条件，$q$ 是 $p$ 的必要条件。联结词优先级通常为 $\neg,\land,\lor,\to,\leftrightarrow$，有歧义时加括号。

\subsection{逻辑等价与基本恒等式}\index{逻辑等价与基本恒等式|hyperpage}
$P\equiv Q$ 表示两个公式在每个赋值下真值相同；$\top,\bot$ 分别表示恒真与恒假。
\begin{gather*}
p\to q\equiv\neg p\lor q\equiv\neg q\to\neg p,\\
p\leftrightarrow q\equiv(p\to q)\land(q\to p)
\equiv(p\land q)\lor(\neg p\land\neg q),\\
p\oplus q\equiv(p\lor q)\land\neg(p\land q),\qquad
\neg\neg p\equiv p,\\
\neg(p\land q)\equiv\neg p\lor\neg q,\qquad
\neg(p\lor q)\equiv\neg p\land\neg q,\\
p\land(q\lor r)\equiv(p\land q)\lor(p\land r),\\
p\lor(q\land r)\equiv(p\lor q)\land(p\lor r),\\
p\lor(p\land q)\equiv p,\qquad p\land(p\lor q)\equiv p,\\
p\land\top\equiv p,\quad p\lor\bot\equiv p,\quad
p\land\neg p\equiv\bot,\quad p\lor\neg p\equiv\top.
\end{gather*}
合取、析取均满足交换律、结合律与幂等律。原命题与逆否命题等价；逆命题 $q\to p$、否命题 $\neg p\to\neg q$ 相互等价，一般不等价于原命题。

\subsection{可满足性、逻辑后承与范式}\index{可满足性、逻辑后承与范式|hyperpage}
公式在某个赋值下为真称可满足，在所有赋值下为真称重言式，在所有赋值下为假称矛盾式。$\Gamma\models Q$ 表示使前提集 $\Gamma$ 全部为真的赋值都使 $Q$ 为真。对有限前提集，
\[
\{P_1,\ldots,P_m\}\models Q
\iff\models(P_1\land\cdots\land P_m)\to Q.
\]
文字是命题变量或其否定。析取范式是若干文字合取式的析取；合取范式是若干文字析取式的合取。消去 $\to,\leftrightarrow$，将否定移至变量前，再用分配律可得到这两种范式。

固定变量次序 $p_1,\ldots,p_n$。对赋值 $a=(a_1,\ldots,a_n)\in\{0,1\}^n$，定义极小项 $m_a$ 和极大项 $M_a$：
\[
m_a=\bigwedge_{i=1}^n\ell_i,\qquad
\ell_i=\begin{cases}p_i,&a_i=1,\\\neg p_i,&a_i=0,\end{cases}
\qquad
M_a=\bigvee_{i=1}^n k_i,\qquad
k_i=\begin{cases}p_i,&a_i=0,\\\neg p_i,&a_i=1.\end{cases}
\]
$m_a$ 仅在赋值 $a$ 下为真，$M_a$ 仅在赋值 $a$ 下为假。因此主析取范式与主合取范式为
\[
P\equiv\bigvee_{a:P(a)=1}m_a,\qquad
P\equiv\bigwedge_{a:P(a)=0}M_a.
\]
空析取约定为 $\bot$，空合取约定为 $\top$。固定变量集后，主范式除项的排列外唯一，普通范式一般不唯一。例如
\[
p\to q\equiv(p\land q)\lor(\neg p\land q)\lor(\neg p\land\neg q)
\equiv\neg p\lor q.
\]

\section{谓词逻辑}
\subsection{量词、辖域与否定}\index{量词、辖域与否定|hyperpage}
默认论域 $D$ 非空。谓词 $P(x)$ 在指定 $x$ 或绑定相应自由变量后才有确定真值；$\forall x\,P(x)$ 表示对所有 $x\in D$ 成立，$\exists x\,P(x)$ 表示至少存在一个成立。受量词约束的出现称约束出现，其余为自由出现。
\begin{gather*}
\neg\forall x\,P(x)\equiv\exists x\,\neg P(x),\qquad
\neg\exists x\,P(x)\equiv\forall x\,\neg P(x),\\
\forall x\in A\,P(x)\equiv\forall x\,(x\in A\to P(x)),\\
\exists x\in A\,P(x)\equiv\exists x\,(x\in A\land P(x)),\\
\exists!x\,P(x)\equiv\exists x\bigl(P(x)\land\forall y\,(P(y)\to y=x)\bigr).
\end{gather*}
受限全称命题在 $A=\varnothing$ 时为真，受限存在命题在此时为假。代换和重命名约束变量时须避免把原来的自由变量误纳入量词辖域。

\subsection{量词交换与分配}\index{量词交换与分配|hyperpage}
同类量词可交换，异类量词一般不可交换：
\begin{gather*}
\forall x\forall y\,R(x,y)\equiv\forall y\forall x\,R(x,y),\qquad
\exists x\exists y\,R(x,y)\equiv\exists y\exists x\,R(x,y),\\
\exists y\forall x\,R(x,y)\ \models\ \forall x\exists y\,R(x,y),\\
\forall x\,(P(x)\land Q(x))\equiv(\forall x\,P(x))\land(\forall x\,Q(x)),\\
\exists x\,(P(x)\lor Q(x))\equiv(\exists x\,P(x))\lor(\exists x\,Q(x)).
\end{gather*}
例如在整数论域上 $\forall x\exists y\,(y>x)$ 为真，而 $\exists y\forall x\,(y>x)$ 为假。全称量词一般不对析取分配，存在量词一般不对合取分配。若 $x$ 不在 $Q$ 中自由出现，则在非空论域上，对 $\circ\in\{\land,\lor\}$ 有
\[
\forall x\,(P(x)\circ Q)\equiv(\forall x\,P(x))\circ Q,\qquad
\exists x\,(P(x)\circ Q)\equiv(\exists x\,P(x))\circ Q.
\]

\section{推理规则与证明方法}
\subsection{常用推理与归结}\index{常用推理与归结|hyperpage}
记 $\Gamma\vdash Q$ 为可以由前提集 $\Gamma$ 按给定推理系统推导出 $Q$。下列命题逻辑规则保持真值：
\begin{gather*}
\frac{p\quad p\to q}{q}\ \text{（假言推理）},\qquad
\frac{\neg q\quad p\to q}{\neg p}\ \text{（拒取式）},\\
\frac{p\to q\quad q\to r}{p\to r}\ \text{（假言三段论）},\qquad
\frac{p\lor q\quad\neg p}{q}\ \text{（析取三段论）},\\
\frac{p\lor q\quad\neg p\lor r}{q\lor r}\ \text{（归结规则）}.
\end{gather*}
合取式可推出各分量；已知两个命题可推出其合取；任一命题可推出包含它的析取。要用归结反证有限前提是否蕴含 $Q$，将前提与 $\neg Q$ 一起化为合取范式，再对其中的子句归结；导出空子句即表明不可满足。

由 $p\to q$ 与 $q$ 推出 $p$ 是肯定后件错误，由 $p\to q$ 与 $\neg p$ 推出 $\neg q$ 是否定前件错误。

\subsection{量词规则与直接、间接证明}\index{量词规则与直接、间接证明|hyperpage}
全称实例化允许由 $\forall x\,P(x)$ 推出 $P(t)$，其中项 $t$ 必须可自由代入 $x$；存在推广允许由 $P(t)$ 推出 $\exists x\,P(x)$。证明全称命题应取任意对象，且不能对它附加前提中没有的特殊条件。

使用 $\exists x\,P(x)$ 时，可引入新的见证符号 $c$ 并假定 $P(c)$；见证不能预先指定，最终结论也不能依赖该新符号。不同存在命题的见证一般不能视为同一对象。
\begin{gather*}
\text{直接证明：假定 }p\text{ 并推出 }q\text{，从而证明 }p\to q;\\
\text{逆否证明：证明 }\neg q\to\neg p\text{，从而证明 }p\to q;\\
\text{反证法：由 }\neg p\text{ 导出矛盾，从而证明 }p.
\end{gather*}
证明等价命题需检查两个方向；否定全称命题只需给出一个反例。

\subsection{数学归纳、强归纳与结构归纳}\index{数学归纳、强归纳与结构归纳|hyperpage}
对整数 $n_0$，数学归纳法为
\[
P(n_0)\ \land\ \forall n\ge n_0\,(P(n)\to P(n+1))
\quad\Longrightarrow\quad\forall n\ge n_0\,P(n).
\]
强归纳法将归纳步骤改为
\[
\forall n\ge n_0\left[\left(\bigwedge_{k=n_0}^{n}P(k)\right)\to P(n+1)\right].
\]
两种方法在自然数上等价，均须验证起始情形；若递推依赖多个初值，应检查足够多的起始情形。结构归纳用于由基本对象经构造规则有限生成的对象：先检查基本对象，再证明每条构造规则保持所需性质。自然数归纳是其特例，更一般的超限归纳见第\ref{chap:set-theory}章。

\section{关系运算与有限表示}
\subsection{逆关系、复合与关系幂}\index{逆关系、复合与关系幂|hyperpage}
设 $R\subseteq X\times Y$、$S\subseteq Y\times Z$，约定先应用 $R$ 再应用 $S$ 的复合为
\begin{gather*}
R^{-1}=\{(y,x):(x,y)\in R\},\\
S\circ R=\{(x,z):\exists y\in Y,\ (x,y)\in R\land(y,z)\in S\},\\
(T\circ S)\circ R=T\circ(S\circ R),\qquad
(S\circ R)^{-1}=R^{-1}\circ S^{-1}.
\end{gather*}
对 $X$ 上的关系，单位关系 $I_X=\{(x,x):x\in X\}$，并定义 $R^0=I_X$、$R^{k+1}=R\circ R^k$；于是 $R^r\circ R^s=R^{r+s}$。$(x,y)\in R^k$ 表示存在长度恰为 $k$ 的关系链，不要求各点不同。
\begin{gather*}
R\text{ 自反}\iff I_X\subseteq R,\qquad
R\text{ 对称}\iff R=R^{-1},\\
R\text{ 反对称}\iff R\cap R^{-1}\subseteq I_X,\qquad
R\text{ 传递}\iff R\circ R\subseteq R.
\end{gather*}
反对称性允许对角元素；它不等同于 $R\cap R^{-1}=\varnothing$。

\subsection{布尔矩阵与关系计算}\index{布尔矩阵与关系计算|hyperpage}
对有限有序集 $X=\{x_1,\ldots,x_m\}$、$Y=\{y_1,\ldots,y_n\}$，关系矩阵定义为 $(M_R)_{ij}=1$ 当且仅当 $(x_i,y_j)\in R$，否则为零。布尔矩阵乘积记为 $\odot$：
\[
(A\odot B)_{ij}=\bigvee_k(a_{ik}\land b_{kj}).
\]
并、交、补关系对应逐元析取、合取和否定，且
\[
M_{R^{-1}}=M_R^{\mathsf T},\qquad
M_{S\circ R}=M_R\odot M_S,\qquad
M_{R^k}=M_R^{\odot k}.
\]
由于矩阵以起点为行、终点为列，复合关系的书写次序与矩阵乘积次序相反。普通整数矩阵乘法计算关系链条数，布尔乘法只记录是否存在；两者不可混淆。

\subsection{关系闭包与Warshall递推}\index{关系闭包与Warshall递推|hyperpage}
包含 $R$ 的最小自反、最小对称和最小传递关系分别为
\[
r(R)=R\cup I_X,\qquad s(R)=R\cup R^{-1},\qquad
R^+=\bigcup_{k\ge1}R^k.
\]
自反传递闭包为 $R^*=\bigcup_{k\ge0}R^k=I_X\cup R^+$。若 $|X|=n\ge1$，则
\[
R^+=\bigcup_{k=1}^{n}R^k,\qquad
R^*=\bigcup_{k=0}^{n-1}R^k.
\]
$R^+$ 中不能普遍把上限减为 $n-1$，因为返回自身的最短正长度环可能长为 $n$。包含 $R$ 的最小等价关系为 $(R\cup R^{-1})^*$。

令 $W^{(0)}=M_R$。Warshall 算法依次允许 $x_1,\ldots,x_k$ 作为中间点：
\[
w_{ij}^{(k)}=w_{ij}^{(k-1)}\lor
\bigl(w_{ik}^{(k-1)}\land w_{kj}^{(k-1)}\bigr),\qquad k=1,\ldots,n.
\]
$W^{(n)}=M_{R^+}$；若初始化为 $M_R\lor I_n$，则得到 $M_{R^*}$。算法时间复杂度为 $O(n^3)$，这里 $\lor$ 对矩阵逐元计算。

\section{格与布尔代数}
\subsection{格、上下确界与对偶}\index{格、上下确界与对偶|hyperpage}
偏序集 $(L,\le)$ 中，若任意两元素都有最大下界与最小上界，则称为格，记
\[
a\land b=\inf\{a,b\},\qquad a\lor b=\sup\{a,b\}.
\]
这里 $\land,\lor$ 为格运算；在命题逻辑中则解释为合取、析取。两种运算均满足交换律、结合律、幂等律及吸收律，且
\[
a\le b\iff a\land b=a\iff a\lor b=b.
\]
有最小元 $0$ 和最大元 $1$ 的格称有界格。任意子集（包括空集）均有上下确界的格称完备格；每个非空有限格均为完备格。对偶原理允许同时交换 $\le$ 与 $\ge$、$\land$ 与 $\lor$、$0$ 与 $1$，从格的恒等式得到其对偶恒等式。

\subsection{分配格、有补格与布尔代数}\index{分配格、有补格与布尔代数|hyperpage}
分配格满足分配律及其对偶式：
\[
a\land(b\lor c)=(a\land b)\lor(a\land c),\qquad
a\lor(b\land c)=(a\lor b)\land(a\lor c).
\]
在有界格中，$b$ 是 $a$ 的补元指 $a\land b=0$ 且 $a\lor b=1$。每个元素都有补元时称有补格；补元在一般有补格中不一定唯一，在有界分配格中若存在则唯一。布尔代数是有补的有界分配格。

\begin{samepage}
在布尔代数中记 $a$ 的补元为 $a'$，有
\begin{gather*}
(a')'=a,\qquad 0'=1,\quad 1'=0,\\
(a\land b)'=a'\lor b',\qquad(a\lor b)'=a'\land b',\\
a\le b\iff b'\le a'.
\end{gather*}
\end{samepage}
幂集 $\mathcal P(U)$ 在包含序下构成布尔代数，其中 $\land=\cap$、$\lor=\cup$、$A'=U\setminus A$、$0=\varnothing$、$1=U$。正整数在整除序下构成格，交与并分别为最大公因数与最小公倍数。

有限布尔代数中，极小非零元素称原子。若原子集合为 $A$，则每个元素唯一表示为它下方原子的并，且
\[
B\cong\mathcal P(A),\qquad |B|=2^{|A|}.
\]
因此有限布尔代数的元素数必为 $2$ 的幂。

\subsection{布尔函数、对偶与逻辑门}\index{布尔函数、对偶与逻辑门|hyperpage}
$n$ 元布尔函数是映射 $f:\{0,1\}^n\to\{0,1\}$，不同函数共有 $2^{2^n}$ 个。其主范式见本章命题逻辑部分；固定其余变量，可按任一变量展开：
\[
f(x_1,\ldots,x_n)
=(x_1\land f(1,x_2,\ldots,x_n))
\lor(\neg x_1\land f(0,x_2,\ldots,x_n)).
\]
对偶函数为 $f^*(x_1,\ldots,x_n)=\neg f(\neg x_1,\ldots,\neg x_n)$，满足 $(f^*)^*=f$；在表达式中同时交换 $\land,\lor$ 与常量 $0,1$ 即得对偶式。

与非、或非分别定义为 $p\uparrow q=\neg(p\land q)$、$p\downarrow q=\neg(p\lor q)$。与非运算单独即可表达所有布尔函数，例如
\[
\neg p=p\uparrow p,\qquad
p\land q=(p\uparrow q)\uparrow(p\uparrow q),\qquad
p\lor q=(p\uparrow p)\uparrow(q\uparrow q).
\]
或非运算也具有同样的完备性。逻辑表达式可据此转换为仅由一种门组成的电路。
\subsection{布尔化简与共识定理}\index{布尔化简与共识定理|hyperpage}
本节省略合取符号，以 $xy=x\land y$、$x'=\neg x$ 表示布尔运算。常用化简式为
\begin{gather*}
x\lor x'y=x\lor y,\qquad x(x'\lor y)=xy,\\
xy\lor xy'=x,\qquad (x\lor y)(x\lor y')=x,\\
xy\lor x'z\lor yz=xy\lor x'z,\\
(x\lor y)(x'\lor z)(y\lor z)=(x\lor y)(x'\lor z).
\end{gather*}
最后两式为共识定理及其对偶式，可删去冗余项。这里析取 $\lor$ 不同于异或 $\oplus$；例如 $x\lor x=x$，而 $x\oplus x=0$。

\subsection{代数标准形与二元Möbius变换}\index{代数标准形与二元Möbius变换|hyperpage}
在 $\mathbb F_2=\{0,1\}$ 上，加法为异或，乘法为合取。任意 $n$ 元布尔函数有唯一的代数标准形
\[
f(x_1,\ldots,x_n)=\bigoplus_{S\subseteq[n]}a_S\prod_{i\in S}x_i,
\qquad a_S\in\mathbb F_2,\qquad[n]=\{1,\ldots,n\}.
\]
空乘积为 $1$，所以 $a_{\varnothing}$ 为常数项。令 $\mathbf1_T$ 为集合 $T$ 的特征向量，则真值与系数的相互转换为
\[
f(\mathbf1_T)=\bigoplus_{S\subseteq T}a_S,\qquad
a_S=\bigoplus_{T\subseteq S}f(\mathbf1_T).
\]
该二元 Möbius 变换是自身的逆变换。计算时，以子集掩码存储 $b_T=f(\mathbf1_T)$；对每个 $i=1,\ldots,n$，对所有含 $i$ 的 $T$ 更新
\[
b_T\leftarrow b_T\oplus b_{T\setminus\{i\}}.
\]
完成后 $b_S=a_S$，需 $O(n2^n)$ 次二元运算。以普通多项式记号书写时，始终按 $2$ 取模，并使用 $x_i^2=x_i$。例如
\[
\neg x=1\oplus x,\qquad
x\lor y=x\oplus y\oplus xy,\qquad
x\to y=1\oplus x\oplus xy.
\]
\section{有限偏序与纠错编码补充}
离散结构的运算参见\cite{mcs,cameron,shoup}。
\subsection{有限偏序的关联代数与Möbius反演}\index{有限偏序的关联代数与Möbius反演|hyperpage}
有限偏序集 $P$ 上定义 $(f*g)(x,z)=\sum_{x\le y\le z}f(x,y)g(y,z)$。区间恒为一的 $\zeta$ 的卷积逆为 $\mu$，递推为
\[
\mu(x,x)=1,\qquad\mu(x,z)=-\sum_{x\le y<z}\mu(x,y).
\]
若 $g(x)=\sum_{y\le x}f(y)$，则 $f(x)=\sum_{y\le x}\mu(y,x)g(y)$。除数偏序给数论 Möbius 反演，子集偏序给容斥。
\subsection{Birkhoff有限分配格表示}\index{Birkhoff有限分配格表示|hyperpage}
有限分配格 $L$ 中，令 $J(L)$ 为非零的并不可约元，$\mathcal I(J(L))$ 为其下闭集格。映射
\[
x\longmapsto\{j\in J(L):j\le x\},\qquad
I\longmapsto\bigvee_{j\in I}j
\]
互逆，给出 $L\cong\mathcal I(J(L))$。下闭集中的交、并对应格的交、并；一般非分配有限格没有此表示。
\subsection{Dilworth与Mirsky分解}\index{Dilworth与Mirsky分解|hyperpage}
有限偏序集的宽度 $w$ 为最大反链大小，高度 $h$ 为最大链大小。则
\[
w=\min\{\text{把 }P\text{ 分割成链的条数}\},\qquad
h=\min\{\text{把 }P\text{ 分割成反链的条数}\}.
\]
第一式为 Dilworth 定理，第二式为其对偶型 Mirsky 定理。这里是互不相交的分割；空偏序集可统一取 $w=h=0$。
\subsection{线性码的距离、校验与球界}\index{线性码的距离、校验与球界|hyperpage}
有限域 $\mathbb F_q$ 上 $[n,k,d]$ 线性码为 $k$ 维子空间 $C$，$d=\min_{c\ne0}\operatorname{wt}(c)$。满秩校验矩阵 $H$ 满足 $C=\ker H$，并有
\[
k\le n-d+1,\qquad
q^k\sum_{j=0}^{t}\binom nj(q-1)^j\le q^n,
\quad t=\left\lfloor\frac{d-1}{2}\right\rfloor.
\]
两式分别为 Singleton 界与 Hamming 球界。综合征 $Hy$ 区分误差的陪集；唯一纠错保证适用于至多 $t$ 个符号错误。
\clearpage
\end{document}
